We need to compute the initial condition for the system for each evaluation of the cost function so that we have a periodic trajectory, 
meaning that the last states of the trajectory are the same as the initial condition.

We can have different dynamics for each switching instant, so the system behaves as 

\begin{align}
    t \in \big[0, t_1\big]   &: \dot{x}= A_{\sigma(0)} x + b_{\sigma(0)}  \nonumber \\
    t \in \big[t_1, t_2\big] &: \dot{x}= A_{\sigma(t_1)} x + b_{\sigma(t_1)} \nonumber \\
    &\qquad \vdots \nonumber \\
    t \in \big[t_{N-1}, t_N\big] &: \dot{x}= A_{\sigma(t_{N-1})} x + b_{\sigma(t_{N-1})}
\end{align}

We can rewrite the solution for the first region as

\noindent
\begin{align}
    x(t_1) &= \green{e^{A_{\sigma(0)} t_1}} x(0) + 
            \orange{\int_{0}^{t_1}{e^{A_{\sigma(0)}(t_1 - \tau)} b_{\sigma(0)} d\tau}} \nonumber \\
    x(t_1) &= \green{F_0} x(0) + \orange{g_0}
\end{align}

\noindent
then the behavior of the other regions of the system is

\noindent
\begin{align}
    x(t_1) &= F_0 x(0) + g_0 \nonumber \\
    x(t_2) &= F_1 x(t_1) + g_1 \nonumber \\
    &\qquad \vdots \nonumber \\ 
    x(t_N) &= F_{N-1} x(t_{N-1}) + g_{N-1}
    \label{cap2-get-x0-multi-equations}
\end{align}

\noindent
with

\noindent
\begin{align}
    F_i = e^{A_{\sigma(t_i)}(t_{i+1} - t_i)} 
\end{align}
\begin{align}
    g_i = \int_{t_i}^{t_{i+1}}{e^{A_{\sigma(t_i)}(t_{i+1} - \tau)} b_{\sigma(t_i)} d\tau} = 
    \int_{0}^{t_{i+1} - t_i}{e^{A_{\sigma(t_i)}\tau} b_{\sigma(t_i)} d\tau}
\end{align}

From \eqref{cap2-get-x0-multi-equations} we can establish the relation between \highlight{$x(t_N)$} and \highlight{$x(0)$}


\begin{align}
    \begin{split}
        x(t_N) &= F_{N-1}F_{N-2} \ldots F_0 x(0) + \\
        &\quad \green{F_{N-1}F_{N-2} \ldots F_1 g_0 + F_{N-1}F_{N-2} \ldots F_2 g_1 + \ldots +} \green{F_{N-1} g_{N-2} + g_{N-1}}
    \end{split} \nonumber \\
    x(t_N) &= F_{N-1} F_{N-2}\ldots F_0 x(0) + \green{c}
\end{align}

We know that \highlight{$x(t_N) = x(0)$}, so the solution for \highlight{$x(0)$} is

\noindent
\begin{align}
    x(0) = F_{N-1}F_{N-2} \ldots F_0 x(0) + c \nonumber \\
    \Big(I - F_{N-1}F_{N-2} \ldots F_0\Big) x(0) = c \nonumber \\
    x(0) = \Big(I - F_{N-1}F_{N-2} \ldots F_0\Big)^{-1} c
    \label{cap2-equacao-matriz-inversa}
\end{align}

\noindent
where
\begin{equation}
    \begin{split}
        c = F_{N-1}F_{N-2} \ldots F_1 g_0 + F_{N-1}F_{N-2} \ldots F_2 g_1 + \ldots + F_{N-1} g_{N-2} + g_{N-1}
    \end{split}
\end{equation}

\noindent
assuming that the inverse indicated in \eqref{cap2-equacao-matriz-inversa} exists.